\begin{example}[\(\mathcal I\)-equivalent, not \(P\)-equivalent]
\label{ex:ch51-cantor}
Let \(X=\mathbf P^1\) with \(\omega=\ddc\log(1+|z|^2)\), normalized by
\(\int_X\omega=1\), and let \(G_p\) be the Green function of
\Cref{sec:ch51-unit-tests}, so that
\[
 G_p(z)=\log\frac{|z-p|^2}{(1+|z|^2)(1+|p|^2)}\leq0,
 \qquad
 \omega+\ddc G_p=\delta_p .
\]
Let \(F\) and \(\mu\) be as in \Cref{lem:ch51-polar-cantor}, and set
\[
 G_\mu(z):=\int_XG_p(z)\,\rd\mu(p).
\]

For \(0\leq s\leq1\), put \(\varphi_s:=sG_\mu\).

\emph{The \(\varphi_s\) are genuine potentials.}  Each \(G_p\leq0\), so
\(G_\mu\leq0\).  For \(z\) at chordal distance at least \(\delta\) from
\(F\) one has \(G_p(z)\geq2\log\delta\) for every \(p\in F\), so
\(G_\mu(z)\geq2\log\delta>-\infty\); in particular
\(G_\mu\not\equiv-\infty\).  Fubini gives
\[
 \omega+\ddc G_\mu=\mu ,
 \qquad
 \omega+\ddc\varphi_s=(1-s)\omega+s\mu\ \geq0,
\]
so \(\varphi_s\in\PSH(X,\omega)\) with \(\varphi_s\leq0\).

\emph{All Lelong numbers vanish.}  On a Riemann surface every
modification is an isomorphism, so the prime divisors over \(X\) are the
points of \(X\), and
\Cref{prop:ch51-lelong-three}, item \cref{itm:ch51-lelong:mass}, gives
\[
 \nu(\varphi_s,x)=(\omega+\ddc\varphi_s)(\{x\})
 =(1-s)\omega(\{x\})+s\mu(\{x\})=0
\]
for every \(x\), because \(\omega\) is smooth and \(\mu\) is nonatomic.
By \Cref{thm:complex-strong-openness}, item
\cref{itm:ch51-so:equivalence}, or directly by Skoda's bound
\Cref{prop:ch51-lelong-calculus}, item \cref{itm:ch51-lc:smooth}, applied
to \(m\varphi_s\) for every \(m\),
\[
 \mathcal I(m\varphi_s)=\cO_X\quad(m\geq1),
 \qquad\text{hence}\qquad
 \varphi_s\sim_{\mathcal I}0 .
\]

\emph{The nonpluripolar mass is \(1-s\).}  The set \(F\) is polar,
hence pluripolar, and \(\mu(F)=1\).  Also \(\omega(F)=0\), since \(F\)
has vanishing Lebesgue measure.  By
\Cref{lem:ch51-npp-bidegree-one} and
\Cref{thm:complex-npp-ledger}, the nonpluripolar product charges no
pluripolar set, so
\[
 \MA_\omega(\varphi_s)
 =\one_{X\setminus F}\bigl((1-s)\omega+s\mu\bigr)
 =(1-s)\omega,
 \qquad
 M_\omega(\varphi_s)=1-s .
\]
The envelope of \(0\) is \(0\), since \(0\) is a competitor and every
competitor is at most zero.  Thus one \(\mathcal I\)-class carries every
mass in \([0,1]\), while
\Cref{cor:ch51-i-canonical-reconstruction} identifies its canonical
normalized \(\mathcal I\)-model representative with \(0\), of mass one.

\emph{Conclusion: \(\varphi_{1/2}\not\sim_P0\).}  Both
\(\varphi_{1/2}\) and \(0\)
have positive mass in the K\"ahler class \(\{\omega\}\), so
\Cref{thm:complex-p-projection} applies to both and
\Cref{thm:ch51-hierarchy-strict}, item \cref{itm:ch51-hs:mass}, gives:
if \(\varphi_{1/2}\sim_P0\), then
\(M_\omega(\varphi_{1/2})=M_\omega(0)\), that is \(\tfrac12=1\).  Hence
\[
 \varphi_{1/2}\sim_{\mathcal I}0,
 \qquad
 \varphi_{1/2}\not\sim_P0,
 \qquad
 \varphi_{1/2}\not\sim0 .
\]
Explicitly, \(P_\omega[\varphi_{1/2}]\) is a model potential with
\(M_\omega(P_\omega[\varphi_{1/2}])=\tfrac12\) and
\(P_\omega[0]=V_\omega=0\) with mass \(1\), so the two envelopes differ.
Since \(\omega+\ddc P_\omega[\varphi_{1/2}]\) has total mass \(1\) and
nonpluripolar mass \(\tfrac12\), it carries mass \(\tfrac12\) on a polar
set; in particular \(P_\omega[\varphi_{1/2}]\) is unbounded, although all its
Lelong numbers vanish.

At the endpoint, \(\varphi_1\) has raw nonpluripolar mass zero but the
same canonical normalized \(\mathcal I\)-model representative \(0\) has
mass one.  Thus neither a raw zero-mass representative nor its raw mass can
be reconstructed from its \(\mathcal I\)-class.

\emph{Direction of the failure.}  One implication does survive:
\(\varphi_{1/2}\leq0\) gives \(\varphi_{1/2}\preceq0\), hence
\(\varphi_{1/2}\preceq_P0\).  What fails is the reverse comparison
\(0\preceq_P\varphi_{1/2}\), even though
\(0\preceq_{\mathcal I}\varphi_{1/2}\) holds.
So the second arrow in \eqref{eq:complex-comparison-hierarchy} is not
reversible, and it is not reversible even after restricting to potentials
of positive mass in a K\"ahler class on a curve.
\end{example}
